\documentclass[../../../../main.tex]{subfiles}
\begin{document}

\begin{flushright}
\footnotesize\textit{【自動文字起こし・要確認】}
\end{flushright}

[解] (i) $\cos(m+n)x = \cos mx \cos nx - \sin mx \sin nx$ (複合同順) の辺々引いて,
$\frac{1}{2} \{ \cos(m-n)x - \cos(m+n)x \} = \sin mx \sin nx$ \qed
したがって,
\[ I_{m,n} = \frac{1}{2} \int_{-\pi}^{\pi} \{ \cos(m-n)x - \cos(m+n)x \} dx \]
\[ = \begin{cases} 0 & (m \neq n) \\ \pi & (m = n) \end{cases} \quad \text{\qed} \]

(ii) 被積分関数を $g_k(x)$ とおく.
\[ g_k(x) = \sin^2 kx - 2(a \sin mx + b \sin nx) \sin kx + (a \sin mx + b \sin nx)^2 \]
\[ = \sin^2 kx + a^2 \sin^2 mx + b^2 \sin^2 nx + 2ab \sin mx \sin nx - 2(a \sin mx + b \sin nx) \sin kx \]
から, (i) より
\[ f(0) = \pi(a^2+b^2) + 2ab I_{m,n} \]
\[ f(k) = \pi(a^2+b^2+1) + 2ab I_{m,n} - 2a I_{m,k} - 2b I_{n,k} \quad (k \ge 1) \]
$5 C_0 = 5 C_5 = 1, 5 C_1 = 5 C_4 = 5, 5 C_2 = 5 C_3 = 10$ だから,
\[ E = \left(\frac{1}{2}\right)^5 \left[ \pi(32a^2 + 32b^2 + 31) + 64ab I_{m,n} - 2a \sum_{k=0}^{5} {}_5C_k I_{m,k} - 2b \sum_{k=0}^{5} {}_5C_k I_{n,k} \right] \]
ここで, $m, n \in \mathbb{N}$ 及び $0 \le k \le 5$ から,
$\sum_{k=0}^{5} {}_5C_k I_{m,k} = \begin{cases} {}_5C_m \cdot \pi & (1 \le m \le 5 \text{ の時}) \\ 0 & (6 \le m \text{ の時}) \end{cases}$
である. 又, $m$ の部を $A$ とおく.

$1^\circ \ m=n \text{ の時}$
$I_{m,n} = \pi$ だから,
$E = \frac{1}{32} [ 32(a+b)^2 + 31 ] \pi \quad (6 \le m)$
$E = \frac{1}{32} [ 32(a+b)^2 - 2(a+b) {}_5C_m + 31 ] \pi \quad (1 \le m \le 5)$
である. $6 \le m$ の時, $a+b=0$ で $\min E = \frac{31}{32}\pi$ をとる.
$1 \le m \le 5$ の時, $a+b = \frac{{}_5C_m}{32}$ で $\min E = \frac{\pi}{32} \left( 31 - \frac{({}_5C_m)^2}{32} \right)$ をとる. ③から
$m=2$ 又は $3$ の時, $\min E$ は最小値 $\frac{223}{256}\pi$ をとる.

$2^\circ \ m \neq n$
$I_{m,n} = 0$ から $A = \pi [32(a^2+b^2)+31]$ である. 対称性から $m < n$ として良い.
$2^5 \cdot E$ の値は以下のようになる.

\begin{center}
\begin{tabular}{|c|c|c|}
\hline
$n \backslash m$ & $1 \le m \le 5$ & $6 \le m$ \\
\hline
$1 \le n \le 5$ & $A - 2\pi(a \cdot {}_5C_m + b \cdot {}_5C_n)$ (⑦) & \\
\hline
$6 \le n$ & $A - 2\pi \cdot a \cdot {}_5C_m$ (ウ) & $A$ (エ) \\
\hline
\end{tabular}
\end{center}

まず, ⑦の時. $a, b \le 0$ ならば $E$ は $a, b$ の単調減少関数だから ($\because {}_5C_k > 0$)
$a=b=0$ で $\min E = \frac{31}{32}\pi$ をとる. $a, b > 0$ の時, $m < n$ だから $(m,n) = (2,3)$ で
$2^5 \frac{1}{\pi} E = 32(a^2+b^2) + 31 - 20(a+b)$
$= 32\left(a - \frac{5}{16}\right)^2 + 32\left(b - \frac{5}{16}\right)^2 + 31 - \frac{25}{4}$
$\ge \frac{99}{4}$ （等号成立は $a=b=\frac{5}{16}$）
から, $a=b=\frac{5}{16}$ で $\min E = \frac{99}{128}\pi$ をとる.

次に(ウ)の時, $a \le 0$ ならば同様に $\min E = \frac{31}{32}\pi$. $a > 0$ の時, $m=2, 3$ で,
$2^5 \frac{1}{\pi} E = 32(a^2+b^2) + 31 - 20a$
$\ge 31 - \frac{25}{8} = \frac{223}{8}$ （等号成立は $a=\frac{5}{16}, b=0$）

(エ)の時は $a=b=0$ で $\min E = \frac{31}{32}\pi$ である.

以上から, $2^\circ$ の⑦の時, 且つ $m \neq n$ で, 対称性から
$(m,n,a,b) = \left(3, 2, \frac{5}{16}, \frac{5}{16}\right), \left(2, 3, \frac{5}{16}, \frac{5}{16}\right)$ \qed

\end{document}
